% ============================================================
% SECTION 10
% SYNTHESIS OF INVARIANCE
% Chapter 02: Invariance
% Book I: Ontology of Viability
% ============================================================

\section{Synthesis of Invariance}
\label{sec:synthesis-of-invariance}

% ------------------------------------------------------------
% 10.1. THE ROLE OF INVARIANCE IN THE ONTOLOGY
% ------------------------------------------------------------

The preceding sections have examined invariance
as a distinct conceptual problem within the
ontology of viability.

The inquiry began with the question of what
remains unchanged when an existent, or one
of its determinations, is considered in relation
to changes and transformations.

The analysis gradually established that
invariance cannot be understood adequately
through an unrestricted assertion of permanence.
It requires the specification of what is being
examined, the transformations under consideration,
and the respect in which something remains
unchanged.

For the purposes of this first edition, invariance
has therefore been provisionally defined as a
relation between an existent, or a determinate
aspect of an existent, and a specified domain
of changes or transformations, such that the
existent or aspect remains unchanged in a
particular respect under those transformations.

This definition provides the working framework
of the present chapter.

It enables the examination of invariance without
prematurely identifying it with identity,
immutability, continuity, structure, or any
other ontological determination.

The purpose of this final section is to consolidate
the distinctions established throughout the
chapter and clarify the conceptual position
that invariance occupies within the broader
development of the ontology of viability.

% ------------------------------------------------------------
% 10.2. THE THREE COMPONENTS OF AN
% INVARIANCE CLAIM
% ------------------------------------------------------------

Every sufficiently specified invariance claim
contains three analytical components.

\begin{enumerate}
    \item \textbf{Subject of invariance:}
    the existent or determinate aspect
    whose invariance is being examined.

    \item \textbf{Domain of transformation:}
    the specified changes or transformations
    in relation to which invariance is asserted.

    \item \textbf{Respect of invariance:}
    the particular determination, feature,
    or relation that remains unchanged.
\end{enumerate}

These components are analytically distinguishable,
but they are not independent ontological entities.
They belong to the specification of a single
invariance claim.

The general structure may be represented
schematically as:

\[
\mathcal{I}(A,\mathcal{D},R)
\]

where:

\begin{itemize}
    \item \(A\) represents the existent or
    determinate aspect under consideration;

    \item \(\mathcal{D}\) represents the
    specified domain of transformations;

    \item \(R\) represents the determination
    or respect in which invariance is asserted.
\end{itemize}

This notation is a conceptual representation
rather than a complete formal definition.
It does not establish a mathematical theory
of invariance or settle its ultimate ontological
status.

Its purpose is to make explicit the conditions
that must be distinguished whenever invariance
is attributed to an existent.

% ------------------------------------------------------------
% 10.3. INVARIANCE AND IDENTITY
% ------------------------------------------------------------

The chapter has consistently distinguished
invariance from identity.

The foundational position of this book is
that being and identity are ontologically
inseparable. Existence is not first established
as an identity-neutral condition to which
identity is subsequently added.

Invariance, however, addresses a different
question.

It concerns whether a specified determination
remains unchanged under a defined domain
of transformations.

The identity of an existent cannot therefore
be reduced to the invariance of one or more
of its determinations.

Likewise, the identification of an invariant
determination does not, by itself, explain
the identity of the existent.

The distinction is essential because the
persistence of a particular feature and
the identity-bearing character of an existent
are not equivalent ontological claims.

The relationship between identity and
transformation remains open for further
development, particularly in relation to
the deeper constitution and dynamics
of being.

% ------------------------------------------------------------
% 10.4. INVARIANCE AND DIFFERENCE
% ------------------------------------------------------------

The analysis of difference established
that difference and invariance are not
necessarily opposed.

Different determinations may coexist
within the same existent, and a transformation
may alter one determination while another
remains unchanged.

Consequently, the existence of difference
does not automatically negate invariance.

At the same time, the same determination
cannot be both changed and unchanged
under the same transformation, in the
same respect, and within the same domain.

This qualification preserves the coherence
of invariance claims.

The distinction between difference and
invariance also prevents the analysis
from assuming that every difference
must produce a loss of identity or that
every invariant determination eliminates
the possibility of difference.

Invariance and difference therefore
identify distinct aspects of the relation
between an existent and its determinations.

% ------------------------------------------------------------
% 10.5. INVARIANCE AND CHANGE
% ------------------------------------------------------------

The relationship between invariance
and change has been established through
the distinction between a specified
unchanged determination and the
alteration of determinations under
transformation.

Invariance is not the absence of all change.

An existent may undergo changes in
some respects while a particular
determination remains invariant
within the relevant domain.

The scope of the invariance claim
must therefore remain limited to
the determination and transformations
actually specified.

Conversely, the identification of
change in one determination does
not establish that every other
determination has changed.

The relationship between invariance
and change is consequently domain-
and determination-dependent.

This formulation prevents both
the reduction of invariance to
absolute permanence and the
reduction of change to the universal
negation of invariance.

% ------------------------------------------------------------
% 10.6. INVARIANCE AND CONTINUITY
% ------------------------------------------------------------

The distinction between invariance
and continuity has shown that
the persistence of a particular
determination does not, by itself,
establish the continuity of an
existent.

Continuity concerns whether and
how an existent continues through
different conditions or phases.

Invariance concerns whether a
specified determination remains
unchanged under a defined domain
of transformations.

An invariant feature may be relevant
to the examination of continuity,
but it cannot automatically serve
as its sufficient explanation.

Likewise, continuity does not
necessarily require the unchanged
persistence of every determination
associated with an existent.

The ontological ground of continuity,
and its relationship to identity
through transformation, remain
open questions.

The present account therefore
preserves the distinction without
assuming that invariance is either
a necessary or a sufficient condition
of continuity.

% ------------------------------------------------------------
% 10.7. INVARIANCE AND STRUCTURE
% ------------------------------------------------------------

The examination of structure
has established that structural
invariance may concern an
organization, a relation, or
a particular structural
determination.

The invariance of one structural
feature does not necessarily
establish the invariance of
the entire structure.

Likewise, the invariance of
an organization under a specified
domain of transformations does
not establish that the structure
is immutable under every possible
transformation.

Structure provides an important
context in which invariance
can be examined, but it must
not be treated as an automatically
established ontological foundation
for identity or continuity.

The deeper ontological status
of structure remains open to
the broader development of
the system.

% ------------------------------------------------------------
% 10.8. FORMS AND DOMAINS OF INVARIANCE
% ------------------------------------------------------------

The chapter has distinguished
several analytical forms of
invariance, including:

\begin{enumerate}
    \item Quantitative invariance.
    \item Qualitative invariance.
    \item Relational invariance.
    \item Structural invariance.
    \item Invariance within temporal
    and non-temporal domains.
    \item Local and general invariance.
\end{enumerate}

These distinctions do not
constitute an exhaustive
classification of invariance.

Nor do they establish mutually
exclusive ontological categories.

Different forms may overlap
within a single case. A particular
determination may be quantitative
and structural, relational and
temporal, or qualitative and local.

The classification is intended
to clarify the different analytical
respects in which invariance
may be examined.

It also reinforces the requirement
that every claim of invariance
must specify its relevant domain
and the determination under
consideration.

% ------------------------------------------------------------
% 10.9. INVARIANCE AND KNOWLEDGE
% ------------------------------------------------------------

The chapter has distinguished
the ontological status of
invariance from the epistemic
conditions through which
it is recognized.

The identification of invariance
involves observation, comparison,
measurement, description,
or reasoning.

These procedures provide
grounds for epistemic claims,
but they do not thereby produce
the ontological relation being
examined.

An existent or one of its
determinations may remain
unchanged without being
recognized by a particular
knower.

Conversely, an apparent
invariance may result from
an incomplete description,
an inadequately specified
domain, or the omission
of relevant determinations.

Knowledge and non-knowledge
are intertwined within the
situated condition of the
knower. The recognition
of one invariant determination
may coexist with non-knowledge
of other aspects that could
be relevant to the analysis.

This does not justify reducing
invariance to knowledge, nor
does it justify ignoring the
epistemic grounds of invariance
claims.

Ontological independence and
epistemic justification must
remain distinct but related
requirements.

% ------------------------------------------------------------
% 10.10. ESTABLISHED COMMITMENTS
% ------------------------------------------------------------

The analysis of this chapter
establishes the following
commitments for the first
edition:

\begin{enumerate}
    \item Invariance is provisionally
    understood as a relation
    between an existent or
    a determinate aspect of
    an existent and a specified
    domain of transformations.

    \item Every invariance claim
    must identify its subject,
    domain of transformation,
    and respect of invariance.

    \item The analytical components
    of an invariance claim
    must not be treated as
    independent ontological
    entities.

    \item Invariance is distinct
    from identity, difference,
    change, continuity, and
    structure, even where
    these concepts are related.

    \item Invariance does not
    imply immutability or
    the absence of all change.

    \item The invariance of
    a particular determination
    does not automatically
    establish the invariance
    of an existent as a whole.

    \item Different forms and
    domains of invariance
    may overlap.

    \item Epistemic recognition
    must be distinguished
    from ontological invariance.

    \item The recognition of
    invariance within a
    limited domain does not
    establish universal
    invariance.

    \item The relational
    definition remains
    provisional and does
    not settle the ultimate
    ontological status of
    invariance.
\end{enumerate}

These commitments delimit
the legitimate use of the
concept within the present
stage of the ontology.

% ------------------------------------------------------------
% 10.11. UNRESOLVED QUESTIONS
% ------------------------------------------------------------

Despite the distinctions
established in this chapter,
several questions remain
unresolved.

First, it remains open
whether invariance has
a deeper ontological
significance beyond its
relational formulation.

Second, the relationship
between invariance and
the identity-bearing
character of existence
has not been fully
determined.

Third, the relationship
between invariance and
the ontological conditions
of continuity requires
further investigation.

Fourth, the status of
invariance across different
realms of being has not
yet been established.

Fifth, the relationship
between invariance and
the dynamics of being
remains open.

Sixth, the analysis has
not determined whether
there are conditions under
which invariance has a
more fundamental role
than the present relational
account can express.

These questions must not
be answered by importing
unexamined assumptions
about immutable substances,
permanent essences, or
fixed ontological foundations.

Their examination requires
the broader framework
that will emerge through
the subsequent development
of the book.

% ------------------------------------------------------------
% 10.11.1. TABLE: SYNTHESIS OF INVARIANCE
% ------------------------------------------------------------

% ============================================================
% TABLE 10 — SYNTHESIS OF INVARIANCE
% Book I: Ontology of Viability
% Part I: Identity and Manifestation
% Chapter 2: Invariance
% ============================================================

\begin{table}[htbp]
    \centering
    \caption{Conceptual Synthesis of Invariance}
    \label{tab:synthesis-of-invariance}

    \renewcommand{\arraystretch}{1.4}

    \begin{tabular}{|p{0.22\textwidth}|p{0.36\textwidth}|p{0.33\textwidth}|}
        \hline
        \textbf{Conceptual Relation} &
        \textbf{Established Distinction} &
        \textbf{Result for the Concept of Invariance} \\
        \hline

        Invariance and identity &
        Identity concerns the existent as identity-bearing; invariance concerns the remaining unchanged of a specified determination under specified transformations. &
        Invariance does not constitute identity or exhaust its meaning. \\
        \hline

        Invariance and difference &
        Difference concerns determinate distinctions, whereas invariance concerns the unchanged status of a particular respect within a specified domain. &
        Difference and invariance may coexist, provided that they concern distinct respects or determinations. \\
        \hline

        Invariance and change &
        Change concerns determinate alteration; invariance concerns the remaining unchanged of a specified determination under a specified transformation. &
        An existent may change in one respect while remaining invariant in another. \\
        \hline

        Invariance and continuity &
        Continuity concerns the continuation of an existent or process; invariance concerns the unchanged status of a particular determination. &
        Invariance does not, by itself, establish continuity, and continuity does not necessarily require a specified invariant feature. \\
        \hline

        Invariance and structure &
        Structure concerns an organized configuration of determinations and relations; structural invariance concerns the preservation of a specified structural determination or relation. &
        The invariance of a structural component does not establish the invariance of the whole structure. \\
        \hline

        Invariance and knowledge &
        Ontological invariance concerns the status of what remains unchanged; epistemic recognition concerns the ways in which a knower identifies and describes that invariance. &
        Knowledge does not constitute ontological invariance, and incomplete knowledge does not establish its absence. \\
        \hline

        Forms and domains of invariance &
        Quantitative, qualitative, relational, and structural forms are analytical distinctions; domains specify the transformations under consideration. &
        These forms are not assumed to be mutually exclusive or to constitute a final classification of ontological kinds. \\
        \hline

        Scope and limits &
        The relational account identifies a subject, a domain of transformations, and a respect in which something remains unchanged. &
        The account provides a working definition without settling the ultimate ontological status of invariance. \\
        \hline

    \end{tabular}
\end{table}

\vspace{0.5em}

\noindent
\textit{General synthesis:}
Invariance is provisionally understood in this chapter as a relation between an existent, or a determinate aspect of an existent, and a specified domain of changes or transformations, such that the existent or aspect remains unchanged in a particular respect under those transformations. This account distinguishes invariance from identity, difference, change, continuity, structure, and epistemic recognition without treating these concepts as wholly unrelated.

% ============================================================
% END OF TABLE 10
% ============================================================

% ------------------------------------------------------------
% 10.12. THE PLACE OF INVARIANCE
% IN THE DEVELOPMENT OF THE SYSTEM
% ------------------------------------------------------------

Invariance occupies a specific
but limited position within
the present development
of the ontology of viability.

It provides a means of examining
what remains unchanged in
particular respects when
an existent or one of its
determinations is considered
under specified transformations.

It therefore contributes
to the conceptual distinction
between identity, difference,
change, continuity, and
structure.

However, invariance does
not independently explain
the existence of an existent,
the origin of its identity,
or the conditions under
which it continues through
transformation.

These questions belong
to the broader ontological
framework and cannot be
settled through the concept
of invariance alone.

The present chapter has
therefore established a
disciplined conceptual
position rather than
a complete metaphysical
theory.

This limitation is deliberate.
The development of the
system requires that
conceptual distinctions
be established before
their deeper ontological
relations are determined.

The question of whether
invariance has a more
fundamental ontological
role may be reconsidered
when the relations among
the realms of being,
light, and the dynamics
of being have been
developed.

% ------------------------------------------------------------
% 10.13. TRANSITION TO THE NEXT CHAPTER
% ------------------------------------------------------------

The examination of invariance
has clarified one dimension
of the relationship between
an existent and the changes
or transformations under
which it is considered.

The next chapter turns
to the problem of unity
and plurality.

This transition is conceptually
necessary because the analysis
of invariance has repeatedly
encountered the coexistence
of different determinations,
relations, and structural
aspects within an existent.

The persistence of one
determination alongside
the alteration of another
raises questions concerning
how plurality is related
to the unity of what is
being examined.

The next chapter will
therefore investigate
the ontological problem
of unity and plurality
without presupposing
that unity must eliminate
difference or that plurality
must dissolve unity.

The distinction between
invariance and the concepts
examined in the present
chapter will remain relevant,
but the central inquiry
will now shift from what
remains unchanged under
transformation to the
relationship between
the one and the many.

% ============================================================
% END OF SECTION 10
% ============================================================